\scanpage{86}
On peut préciser, en introduisant l'Algèbre engendrée par \deleted{$A$} $G(\mathbb Q)$ dans $A$, soit $B$.

On sait que $B=\prod L_i$, chaque $L_i$ ext. quadratique d'un corps de nbres $K_i$ tot. réel. ($K_i$ déterminé uniqmt \deleted{par la condition} \inserted{dans $L_i$} comme l'ens des élémts tot. réels de $L_i$), l'involution est donnée sur $L_i$ par la conjugaison p.r. à $K_i$ (i.e. la conjugaison complexe ordinaire, si $L_i\subset\mathbb C$), la trace est proportionnelle, sur chaque $L_i$, à la trace ordinaire.

Ainsi $B^*$ n'est autre que $H(\mathbb Q)$, où

\[
H=\prod_i\left(\prod_{L_i/\mathbb Q}\mathrm{Gm}_{L_i}\right)
\]

et le groupe unitaire $B^{*u}$ n'est autre que le \deleted{groupe} des pts rationnels de

\[
H^u=\prod_i\left(\prod_{K_i/\mathbb Q}\Gamma_i\right)
\]

où $\Gamma_i$ est la forme tordue de $\mathrm{Gm}$ sur $K_i$ définie par l'extension quadratique $L_i$ de $K_i$.

\deleted{[} \deleted{Ainsi} \unsure{Cel|Cela|Ceci} met en évidence que $H^u$ est un tore ``compact'', \deleted{\unsure{donc il est}} (de nature d'ailleurs assez particulière !), et aussi $U^0$ comme sous-tore de ce dernier\ldots.

NB. Il convient de déterminer $\varepsilon:G\longrightarrow\mathrm{Gm}$ en termes \inserted{de} \deleted{des} \unsure{leurs|les|tores} normes sur les $L_i$, et $j:\mathrm{Gm}\longrightarrow G$ \margin{$=\ldots$}
% archive sheet number: 85

